\clearpage\thispagestyle{empty}
\pdfbookmark[0]{Quellen zur Vesta-Notiz}{vesta1807-sources}
\begin{center}\textbf{Quellen zur Vesta-Notiz}\end{center}
\noindent Diese KI-Übertragung in heutiges Deutsch folgt der Notiz \emph{Vesta} in \href{https://gdz.sub.uni-goettingen.de/id/PPN236007467}{Gauß’ \textit{Werken}, Band~VI} (Göttingen, 1874), S.~282. Der Auszug stammt aus Franz Xaver von Zachs Bericht über Wilhelm Olbers’ Entdeckung der Vesta in der \textit{Monatlichen Correspondenz}, Band~XV, Mai 1807, Beitrag XXXV, S.~502--508. Der hier wiedergegebene Wortwechsel steht auf \href{https://zs.thulb.uni-jena.de/rsc/viewer/jportal_derivate_00237836/Monatliche_Correspondenz_130168688_15_1807_0514.tif}{S.~507}. Von Zach zitiert Olbers’ Brief vom 22. April, in dem dieser wiederum Gauß’ Antwort zur Namenswahl zitiert. Olbers beschreibt auch das von Gauß gewählte Zeichen.
\par\medskip
Gauß’ Brief an Olbers vom 14. April 1807 enthält die Stelle zur Namenswahl in etwas anderer Form und erörtert mehrere Entwürfe für das Zeichen. Er steht in Carl Schillings Ausgabe des Briefwechsels (Berlin, 1900), erste Abteilung, \href{https://archive.org/details/p1wilhelmolberss02olbeuoft/page/n351/mode/1up}{S.~338}--\href{https://archive.org/details/p1wilhelmolberss02olbeuoft/page/n352/mode/1up}{339}. Das editierbare Zeichen folgt der in der Werkausgabe gedruckten Kontur. Das Zeitschriftendigitalisat stammt von der Thüringer Universitäts- und Landesbibliothek Jena. Übertragung und Satz wurden mit KI erstellt.
\clearpage
